\section{Extrema in Several Variables}

\subsection{Critical points}
$\vc x_0$ is a \emph{critical point} of $f$ if $\grad f(\vc x_0) = \vc 0$ or some partial does
not exist. Local extrema of a differentiable function occur only at critical points.

\subsection{Second derivative test (two variables)}
At a critical point $(a,b)$ of $f(x,y)$, with continuous second partials, let
\[ H = \det \begin{pmatrix} f_{xx} & f_{xy} \\ f_{xy} & f_{yy} \end{pmatrix}_{(a,b)}
= f_{xx} f_{yy} - f_{xy}^2. \]
\begin{itemize}[leftmargin=*,nosep]
  \item $H > 0$ and $f_{xx} > 0$: local \textbf{minimum}.
  \item $H > 0$ and $f_{xx} < 0$: local \textbf{maximum}.
  \item $H < 0$: \textbf{saddle point}.
  \item $H = 0$: inconclusive.
\end{itemize}

\begin{example}
$f(x,y) = x^3 - 3xy + y^3$. $\grad f = (3x^2 - 3y, -3x + 3y^2) = 0$ at $(0,0)$ and $(1,1)$.
$f_{xx} = 6x$, $f_{yy} = 6y$, $f_{xy} = -3$, so $H = 36xy - 9$.
At $(0,0)$: $H = -9 < 0$, saddle. At $(1,1)$: $H = 27 > 0$ and $f_{xx} = 6 > 0$, local minimum.
\end{example}

\subsection{Absolute extrema on a closed bounded region}
\begin{enumerate}[leftmargin=*,nosep]
  \item Find critical points in the interior.
  \item Find extrema on the boundary (parametrize the boundary or use Lagrange).
  \item Evaluate $f$ at all candidates; the largest/smallest are absolute extrema.
\end{enumerate}

\subsection{Lagrange multipliers}
To extremize $f(\vc x)$ subject to $g(\vc x) = c$ (with $\grad g \neq 0$):
\[ \grad f = \lambda \grad g, \qquad g(\vc x) = c. \]
Solve the system for $\vc x$ and the multiplier $\lambda$.

\noindent\textbf{Multiple constraints} $g_1 = c_1, \ldots, g_k = c_k$:
$\grad f = \sum_{i=1}^k \lambda_i \grad g_i$.

\begin{example}
Maximize $f(x,y) = xy$ subject to $x^2 + y^2 = 1$.
$\grad f = (y, x)$, $\grad g = (2x, 2y)$, so $y = 2 \lambda x$, $x = 2 \lambda y$. Combined with
$x^2 + y^2 = 1$: $x = \pm y = \pm 1/\sqrt 2$. Maximum value $1/2$ at $(\tfrac{1}{\sqrt 2}, \tfrac{1}{\sqrt 2})$.
\end{example}

\begin{intuition}
At an extremum on the constraint surface, the level set of $f$ is tangent to the constraint
surface; their normals $\grad f$ and $\grad g$ must be parallel.
\end{intuition}

\begin{pitfallbox}
\begin{itemize}[leftmargin=*,nosep]
  \item Lagrange gives \emph{candidates}; compare values to identify the actual max/min.
  \item Check the regularity condition $\grad g \neq 0$ on the constraint set;
        violations require separate consideration.
\end{itemize}
\end{pitfallbox}
